\documentclass[12pt,a4paper]{article}

\usepackage{iftex}

\ifPDFTeX
  \usepackage{cmap}
  \usepackage[T2A]{fontenc}
  \usepackage[utf8]{inputenc}
  \usepackage{lmodern}
  \usepackage[english,russian]{babel}
\else
  \usepackage{fontspec}
  \defaultfontfeatures{Ligatures=TeX}
  \setmainfont{Times New Roman}
  \setsansfont{Arial}
  \setmonofont{Consolas}
  \usepackage{polyglossia}
  \setmainlanguage{russian}
  \setotherlanguage{english}
\fi

\usepackage{amsmath,amssymb,amsthm,mathtools}
\usepackage[a4paper,margin=2.2cm,top=2.4cm,bottom=2.2cm]{geometry}
\usepackage{enumitem}
\usepackage{fancyhdr}
\usepackage[hidelinks]{hyperref}

\providecommand{\SheetNo}{}
\newcommand{\sheetauthor}{}

\newcommand{\sheettitle}[1]{
  \begin{center}
    {\Large\bfseries Серия \SheetNo. #1}\\[0.4em]
    \rule{\linewidth}{0.4pt}
  \end{center}
  \vspace{0.5em}
  \markboth{#1}{#1}
}

\setlength{\headheight}{15pt}
\pagestyle{fancy}
\fancyhf{}
\fancyhead[L]{\small Серия \SheetNo}
\fancyhead[R]{\small\leftmark}
\fancyfoot[C]{\thepage}
\renewcommand{\headrulewidth}{0.4pt}

\newlist{problems}{enumerate}{2}
\setlist[problems]{label=\arabic*., ref=\arabic*,
                   leftmargin=*, itemsep=0.8em, topsep=0.8em}

\newcommand{\cyralphnum}[1]{
  \ifcase#1\or а\or б\or в\or г\or д\or е\or ж\or з\or и\or к\or л\or м\or
  н\or п\or р\or с\or т\or у\or ф\or х\or ц\or ч\or ш\or щ\or э\or ю\or я\fi}
\newcommand{\cyralph}[1]{\cyralphnum{\value{#1}}}
\AddEnumerateCounter{\cyralph}{\cyralphnum}{щ}

\newlist{parts}{enumerate}{2}
\setlist[parts]{label=\cyralph*), ref=\cyralph*,
                leftmargin=1.5em, itemsep=0.25em, topsep=0.25em}

\newcommand{\N}{\mathbb{N}}
\newcommand{\Z}{\mathbb{Z}}
\newcommand{\Q}{\mathbb{Q}}
\newcommand{\R}{\mathbb{R}}
\providecommand{\C}{}
\renewcommand{\C}{\mathbb{C}}
\providecommand{\NOD}{}
\renewcommand{\NOD}{\operatorname{\text{НОД}}}
\providecommand{\NOK}{}
\renewcommand{\NOK}{\operatorname{\text{НОК}}}


\def\SheetNo{2}

\begin{document}
\sheettitle{Множества, функции, отношения.}
\begin{problems}


\item Пусть $A_1,A_2,\ldots,A_n$ "--- конечные множества. Докажите, что
\[
\Bigl|\,\bigcup_{i=1}^{n}A_i\Bigr|
=\sum_{i}|A_i|-\sum_{i<j}|A_i\cap A_j|+\sum_{i<j<k}|A_i\cap A_j\cap A_k|-\ldots+(-1)^{n-1}|A_1\cap A_2\cap\ldots\cap A_n|,
\]
то есть
\[
\Bigl|\,\bigcup_{i=1}^{n}A_i\Bigr|
=\sum_{\varnothing\neq I\subseteq\{1,\ldots,n\}}(-1)^{|I|-1}\Bigl|\,\bigcap_{i\in I}A_i\Bigr|.
\]

\item 
Пусть $M$ "--- множество всех трехэлементных подмножеств натурального ряда. Рассмотрим отображение $f\colon M \to \mathbb{N}$, которое каждому трехэлементному множеству ставит в соответствие его второй по величине элемент. Является ли это отображение а)~инъекцией б)~сюръекцией?

\item Постройте пример:
\begin{enumerate}[label=\alph*)]
  \item рефлексивного, симметричного, но не транзитивного отношения;
  \item рефлексивного, антисимметричного, но не транзитивного отношения;
  \item отношения, которое не симметрично и не антисимметрично;
  \item отношения, которое симметрично и антисимметрично.
\end{enumerate}

\item Рассмотрим на множестве $\mathbb{R}$ бинарное отношение $R$. Чему равно $R\circ R$, если $R(x,y)$ означает:
\begin{enumerate}[label=\alph*)]
  \item $y=x+1$;
  \item $x+y=1$;
  \item $xy=1$;
  \item $xy>1$.
\end{enumerate}

\item Является ли отношение $R$, такое что $(a,b)R(c,d)$, если $a+d=b+c$ на $\mathbb{N}\times\mathbb{N}$ отношением эквивалентности?

\item Композиция отношений $R$ и $S$ это отношение $T=RS$, где $xTy$, если найдется $z$, такой что $xRz$ и $zSy$. Пусть $R$ и $S$ "--- транзитивные отношения на $A$. Будет ли транзитивной их композиция?

\item Дано множество $X$, $|X|=n$. Сколько существует:
\begin{enumerate}
\item  бинарных отношений на множестве $X$?
\item рефлексивных?
\item симметричных?
\item антисимметричных?
\end{enumerate}

\item Какие из равенств верны для любых множеств $A,B,C$?
\begin{enumerate}
  \item $(A\cap B)\cup C=(A\cup C)\cap(B\cup C)$;
  \item $(A\cup B)\cap C=(A\cap C)\cup(B\cap C)$;
  \item $(A\cup B)\setminus C=(A\setminus C)\cup B$;
  \item $(A\cap B)\setminus C=(A\setminus C)\cap B$;
  \item $A\setminus(B\cup C)=(A\setminus B)\cap(A\setminus C)$;
  \item $A\setminus(B\cap C)=(A\setminus B)\cup(A\setminus C)$.
\end{enumerate}

\item Проведите подробное доказательство верных равенств предыдущей задачи, исходя из определений. (Докажем, что множества в левой и правой частях равны. Пусть $x$ "--- любой элемент левой части равенства. Тогда\ldots{} Поэтому $x$ входит в правую часть. Обратно, пусть\ldots) Приведите контрпримеры к неверным равенствам.

\item На какое наименьшее число групп можно разбить все пары элементов данного $n$-элементного множества так, чтобы любые две пары из одной группы имели общий элемент?

\item Докажите, что в любой выборке из 52 положительных целых чисел найдутся хотя бы два, у которых либо их сумма, либо их разность делится на 100.
 
\item Пусть $M$ "--- количество всевозможных способов расставить слонов на шахматной доске так, чтобы никакие два из них не били друг друга. (Количество слонов может быть любым натуральным числом, даже 1.) Докажите, что $M+1$ "--- квадрат некоторого натурального числа.
 
\item У конечного множества $A$ выбрано несколько подмножеств. Каждое подмножество состоит хотя бы из двух элементов, а любые два пересекающихся подмножества имеют хотя бы два общих элемента. Докажите, что элементы множества $A$ можно покрасить в два цвета так, чтобы каждое из выбранных подмножеств содержало элементы обоих цветов.

\item Сколько различных выражений для множеств можно составить из переменных $A$ и $B$ с помощью (многократно используемых) операций пересечения, объединения и разности? (Два выражения считаются одинаковыми, если они равны при любых значениях переменных.) Тот же вопрос для трёх множеств и для $n$ множеств. (Ответ в общем случае: $2^{2^n-1}$.)

\item Докажите, что $f^{-1}\bigl(g^{-1}(Z)\bigr)=(g\circ f)^{-1}(Z)$ для любого $Z\subseteq C$ и любых функций $f$ из множества $A$ в множество $B$ и $g$ из множества $B$ в множество $C$. Здесь $f^{-1}$ обозначает полный прообраз множества.

\item На плоскости нарисовано некоторое количество равносторонних треугольников. Они не пересекаются, но могут иметь общие участки сторон. Мы хотим покрасить каждый треугольник в какой-нибудь цвет так, чтобы те из них, которые соприкасаются, были покрашены в разные цвета (треугольники, имеющие одну общую точку, могут быть покрашены в один цвет). Хватит ли для такой раскраски двух цветов?

\item Из натуральных чисел от 1 до $2n$ выбрали $n+1$ число. Докажите, что одно из этих чисел делится на другое.

\item Верно ли для любых $f$, $Y_1$, $Y_2$ равенство $f^{-1}(Y_1\cap Y_2)=f^{-1}(Y_1)\cap f^{-1}(Y_2)$?
 
\item Пусть $f$ всюду определена, а $g$ не всюду определена. Верно ли, что композиция $g\circ f$ не всюду определена?
 
\item Докажите, что для любой функции $f\colon A\to A$, отображающей конечное множество $A$ в себя, найдутся различные целые положительные $m$, $n$, при которых $f^m$ и $f^n$ совпадают.
 
\item Мы уже определили $f^n$ при целом $n\geqslant 2$, а также $f^{-1}$ для биекций. Как надо определить $f^0$, $f^1$ и отрицательные степени (для биекций), чтобы выполнялось равенство
\[
f^{m+n}=f^m\circ f^n
\]
(при любых целых $m$ и $n$)?

 \item Найдите количество:
\begin{enumerate}
  \item[a)] неубывающих функций $f\colon\{1,2,\ldots,n\}\to\{1,2,\ldots,m\}$;
  \item[b)] неубывающих инъекций $f\colon\{1,2,\ldots,n\}\to\{1,2,\ldots,m\}$;
  \item[c)] неубывающих сюръекций $f\colon\{1,2,\ldots,n\}\to\{1,2,\ldots,m\}$.
\end{enumerate}
Функция неубывающая, если $x\leqslant y$ влечёт $f(x)\leqslant f(y)$.

\item Докажите, что среди чисел Фибоначчи найдется число, оканчивающееся тремя нулями.

\item В школе действует несколько кружков. Среди них есть кружок по топологии, который никто не посещает, и кружок по обществознанию, на который ходят все 2020 учеников школы. Списочные составы любых двух кружков различны. Кроме того, для любых двух кружков $A$ и $B$
\begin{enumerate}[label=\arabic*)]
  \item найдётся кружок $C$, который посещают в точности те ученики, которые посещают как $A$, так и $B$;
  \item найдётся кружок $D$, который посещают в точности те ученики, которые посещают хотя бы один из кружков $A$ или $B$.
\end{enumerate}
Докажите, что какой-то ученик посещает не менее половины всех кружков.

\item Дано натуральное число $n$, и множество $S=\{1,2,\ldots,n\}$. Найдите наименьшее значение величины $|S\triangle A|+|S\triangle B|+|S\triangle C|$, где $A$ и $B$ "--- непустые конечные множества натуральных чисел и $C=A+B=\{a+b \mid a\in A,\, b\in B\}$.

\end{problems}
\end{document}
